\section{Quantitative penalties for genuine Markoff sources}
\label{sec:markoff-application}

The equality graph describes paths with total excess zero. Near-minimizers
can also visit zero-slack components that have a positive cost of entry
or exit. To understand them, we must examine the entire zero-slack graph,
not only the six states surviving the equality test.

\subsection{The full critical set}

We first record the structure of the zero-slack subgraph.

\begin{table}[tb]
\centering
\caption{Every cyclic zero-slack component at norm $89$.
The frequency $\theta_i$ counts paired letters equal to $2$.
The oscillation $D_i$ is that of the rational count potential
$g_i=H_i/d_i$. Every internal edge satisfies
$\mathbf1_{\{j=2\}}-\theta_i=g_i(S')-g_i(S)$.}
\label{tab:critical-components}
\begin{tabular}{rrrr}
\toprule
Number of components & States per component & $\theta_i$ & $D_i$\\
\midrule
1 & 6 & $1/5$ & $1$\\
2 & 3 & $1/3$ & $2/3$\\
2 & 1 & $0$ & $0$\\
2 & 1 & $1$ & $0$\\
\bottomrule
\end{tabular}
\end{table}

\begin{proposition}[The critical components at norm $89$]
\label{prop:critical-components}
The zero-slack subgraph of \eqref{eq:potential-inequalities} has $12684$
edges and $11411$ strongly connected components, of which exactly seven
are cyclic. Their frequencies and oscillations are those listed in
\cref{tab:critical-components}, and every internal edge of each satisfies
the count-potential identity \eqref{eq:count-potential}. The longest path
in the condensation, weighted as in \eqref{eq:condensation-budget}, has
\begin{equation}
 L=58,\qquad K=59.
 \label{eq:axis-constants}
\end{equation}
\end{proposition}

\begin{proof}
All of this concerns the finite graph of \cref{prop:closed-graph}, so it
is again a finite assertion about integers, and we verified it by
computer in exact arithmetic, twice and independently. A path realizing
the maximum traverses $56$ condensation edges together with the
six-state component, whose vertex weight is two; this accounts for the
value of $L$. The states and internal edges of the seven cyclic
components are listed in the supplementary material.
\end{proof}

\begin{theorem}[Norm-$89$ balanced stability]
\label{thm:axis-stability}
Under the hypotheses of \cref{thm:endpoint-89}, suppose the core $w$ is
balanced. With $n=|w|$, $m=|w|_2$ and
$\Theta_{89}=\{0,1/5,1/3,1\}$,
\begin{equation}
 \min_{\theta\in\Theta_{89}}|m-\theta n|
 <59(\Delta+1)+2.
 \label{eq:axis-core-bound}
\end{equation}
If the source is a genuine reduced occurrence of slope $P/Q$, then
\begin{equation}
 \min\{P,\,Q-P,\,|P-Q/5|,\,|P-Q/3|\}\leqslant118\Delta+121.
 \label{eq:axis-farey-bound}
\end{equation}
\end{theorem}

\begin{proof}
\Cref{prop:closed-graph} supplies the nonnegative integer weights
required by \cref{thm:balanced-stability}, and
\cref{prop:critical-components} gives $L=58$ and $K=59$. At most two paired
letters are consumed in the sign-resolution prefix, so the theorem with
$h=2$ gives \eqref{eq:axis-core-bound}.

For a genuine occurrence the auxiliary slope can be dispensed with. By
\cref{cor:farey-compatible} the Farey slope $P/Q$ is itself compatible with
the core, so \cref{lem:budget} may be applied to $\beta=P/Q$. Its hypothesis
$\delta\leqslant1$ holds with room to spare: the complementary intervals of
$\Theta_{89}$ in $[0,1]$ are $[0,1/5]$, $[1/5,1/3]$ and $[1/3,1]$, whose
half-lengths are $1/10$, $1/15$ and $1/3$, so
\begin{equation}
 \delta=\operatorname{dist}(P/Q,\Theta_{89})\leqslant\tfrac13 .
 \label{eq:delta-third}
\end{equation}
The left side of \eqref{eq:axis-farey-bound} is exactly $Q\delta$. Writing
$\mathbf1_{\ne}$ for the indicator of $\varepsilon\ne\varnothing$, the count
dictionary \eqref{eq:source-counts} gives $Q=2n+2+\mathbf1_{\ne}$, whence
$$
 Q\delta=2n\delta+(2+\mathbf1_{\ne})\delta
 \leqslant2\bigl(h+\Delta+(\Delta+1)L\bigr)+3\delta .
$$
With $h=2$, $L=58$ and \eqref{eq:delta-third} the right side is
$2\bigl(2+\Delta+58(\Delta+1)\bigr)+1=2(59\Delta+60)+1=118\Delta+121$.
\end{proof}

The gain over the route through \eqref{eq:axis-core-bound}, which would
give $118\Delta+124$, is small. What matters is that no slope other than
the occurrence's own enters the argument: \cref{cor:farey-compatible}
replaces an existence statement about balanced words by an identity about
the Farey address, and this is the only place where the narrowness recorded
in \cref{prop:thin-language} is actually used.

\begin{corollary}[Slope-dependent linear penalty]
\label{cor:farey-penalty}
If a genuine source in \cref{thm:axis-stability} satisfies
$\operatorname{dist}(P/Q,\Theta_{89})\geqslant\eta>0$, then
$$
 \Delta\geqslant\frac{\eta Q-121}{118}.
$$
Thus any sequence of genuine sources with bounded norm-$89$ excess and
unbounded denominators can accumulate only at the four critical slopes.
\end{corollary}

\begin{proof}
For each $\theta\in\Theta_{89}$ one has
$|P-\theta Q|=Q\,|P/Q-\theta|\geqslant\eta Q$, so the left side of
\eqref{eq:axis-farey-bound} is at least $\eta Q$; rearranging that
inequality gives the bound. The last assertion is the contrapositive:
along a sequence with $\Delta$ bounded and $Q\to\infty$ the quantity
$(\eta Q-121)/118$ is eventually larger than any bound on $\Delta$, so no
positive $\eta$ can persist.
\end{proof}

The constants arise from the full condensation and were not optimized
using the genuine source language. In particular, the last assertion is
a necessary condition on accumulation points, not an existence theorem
at each critical slope.

\subsection{An infinite arithmetic family with growing excess}

\begin{theorem}[A genuine infinite family]
\label{thm:infinite-family}
There is an infinite sequence of genuine reduced Markoff occurrences
$P_j/Q_j$ whose labels have exact $89$-adic valuation one, for which
$P_j,Q_j\longrightarrow\infty$ and
\begin{equation}
 \frac{P_j}{Q_j}\longrightarrow
 \rho=\frac{1+\Phi}{4+5\Phi},\qquad
 \Phi=\frac{1+\sqrt5}{2}\quad\text{the golden ratio}.
 \label{eq:family-limit}
\end{equation}
For all sufficiently large $j$, their complete norm-$89$ excesses satisfy
\begin{equation}
 \Delta_j\geqslant\frac{Q_j}{1180(4+5\Phi)}-\frac{121}{118}.
 \label{eq:family-penalty}
\end{equation}
\end{theorem}

\begin{proof}
Start with the Markoff triple $(1,34,89)$ and iterate
\begin{equation}
 F(a,b,c)=(b,c,3bc-a).
 \label{eq:vieta-forward}
\end{equation}
The identity $a^2+b^2+c^2=3abc$ is preserved. If $1\leqslant a<b<c$,
then $3bc-a>c$, so the positive triples continue forward in the Markoff
tree and their maxima increase strictly.

Modulo $89^2$, the same map is a permutation because it has the polynomial
inverse
$$
 F^{-1}(a,b,c)=(3ab-c,a,b).
$$
The orbit of the initial residue triple is therefore periodic from its
first term. Let $T>0$ be a period. For every $j\geqslant0$, the last
coordinate of $F^{jT}(1,34,89)$ is congruent to $89$ modulo $89^2$.
It consequently has valuation exactly one at $89$.

The initial labels $34$ and $89$ occur at the adjacent Farey vertices
$1/4$ and $1/5$; their other common neighbour $0/1$ has label one. The
Vieta step is exactly the Farey mediant: the Christoffel word of the
mediant of two Farey neighbours is the concatenation of their Christoffel
words, so $\mu$ of the mediant is the product of the two matrices, and
the Fricke identity $\tr(XY)+\tr(XY^{-1})=\tr X\,\tr Y$ turns this into
$$
 m_{\text{mediant}}=3m_\alpha m_\beta-m_{\text{other}}
$$
for the three labels concerned. Every step of \eqref{eq:vieta-forward}
therefore replaces the older vertex by the mediant of the two retained
vertices. Hence the two Farey columns follow
Fibonacci addition from $(1,4)$ and $(1,5)$. Their determinant has absolute
value one, so all resulting fractions are reduced. Both coordinates grow
without bound, and their ratios tend to \eqref{eq:family-limit}. The
subsequence at multiples of $T$ has the same limit.

This limit lies strictly between $1/5$ and $1/4$. Its nearest critical
slope is $1/5$, and
$$
 \rho-\frac15=\frac{1}{5(4+5\Phi)}>0.
$$
For sufficiently large $j$, the distance of $P_j/Q_j$ from $\Theta_{89}$
is at least half this positive number. Apply \cref{cor:farey-penalty}
with $\eta=1/[10(4+5\Phi)]$ to obtain
\eqref{eq:family-penalty}.
\end{proof}

The argument needs neither the value of $T$ nor a conjecture about prime
values of recurrence sequences. It is explicit in the sense that the
family is described by a rule and not merely shown to exist; the period
$T$ is large, so the members beyond the first are not small numbers. Its modular periodicity selects the
required valuation, while the real Farey dynamics places the source
slopes away from the critical set. These two elementary mechanisms
make the quantitative theorem nonvacuous on genuine sources with
unbounded numerators.

\subsection{Why the equality slope alone is insufficient}

The six-state equality graph has frequency $1/5$, but the full critical
graph also contains the state
$$
 S=\begin{pmatrix}101&19\\19&82\end{pmatrix}
  =\begin{pmatrix}10&1\\1&9\end{pmatrix}^{\!2}.
$$
This state commutes with $M(1)^2$. Exactly two forced row subtractions
return it after one light paired letter, so the loop has zero cost.
A fixed entry and primitive exit give the explicit family
\begin{equation}
 q_j=\chi(21221122112\,1^j\,2122121)2,\qquad j\geqslant0.
 \label{eq:zero-light-family}
\end{equation}
The prefix enters $S$ with cumulative weight $15$. The suffix reaches
$(109,49;43,92)$ with additional weight $20$; its terminal vector is
$89(3,2)^{\mathsf T}$ and its terminal charge is $1$.
Thus every member has exact norm valuation one at $89$ and completed
excess $36$. Its inferred counts are $P=21$ and $Q=38+2j$.
This rules out a cost bound proportional only to $|Q-5P|$ on the full
regular language.

Every displayed core contains both $11$ and $22$, so it is unbalanced.
The family is consequently not a family of genuine Markoff sources and
does not contradict \cref{thm:axis-stability}. It explains why the
zero-slack graph outside the exact equality language cannot simply be
discarded when studying near-minimizers.
